\chapter{Chapter07 — CSP Derandomization}
%%% generated by the blueprint pipeline from FABL.Chapter07.CSPDerandomization
\section{Overview}

Declarations of the Lean module \texttt{FABL.Chapter07.CSPDerandomization}.

\section{Declarations}

\begin{theorem}[Uniform binary affine equation]
\lean{FABL.expect_f₂_add_eq}
\leanok
\uses{FABL.𝔽₂}
\difficulty{4}
Let $a,r\in\mathbb F_2$. For a uniformly random $h\in\mathbb F_2$, the probability
that $h+a=r$ is $1/2$.
\end{theorem}

\begin{theorem}[Acceptance probability of a nonzero linear predicate]
\lean{FABL.expect_maxLinearPredicate}
\label{FABL.expect_maxLinearPredicate}
\leanok
\uses{FABL.MaxLinearPredicate, FABL.expect_f₂_add_eq, FABL.maxLinearTemplate, FABL.𝔽₂}
\difficulty{5}
Let $r$ be a nonnegative integer, and let a predicate specify a linear equation over
$\mathbb F_2$ in $k$ variables, where $1\le k\le r$, with any right-hand side in
$\mathbb F_2$. For uniformly random values of those variables, the equation holds with
probability $1/2$.
\end{theorem}

\begin{definition}[Assignment split by constraint scope]
\lean{FABL.maxLinearScopeSplitEquiv}
\label{FABL.maxLinearScopeSplitEquiv}
\leanok
\uses{FABL.CSPConstraint, FABL.maxLinearTemplate, FABL.𝔽₂}
Let $V$ be a finite set, let $r$ be a natural number, and let $C$ be a constraint from
the maximum linear template of parameter $r$ on $V$. There is an equivalence between
assignments $V \to \mathbb{F}_2$ and pairs consisting of an assignment of values in
$\mathbb{F}_2$ to the ordered variables in the scope of $C$, and an assignment of values
in $\mathbb{F}_2$ to the variables of $V$ outside that scope.
\end{definition}

\begin{theorem}[First coordinate of the scope decomposition]
\lean{FABL.maxLinearScopeSplitEquiv_fst_apply}
\label{FABL.maxLinearScopeSplitEquiv_fst_apply}
\leanok
\uses{FABL.CSPConstraint, FABL.maxLinearScopeSplitEquiv, FABL.maxLinearTemplate, FABL.𝔽₂}
\difficulty{3}
Let $V$ be a finite set, let $r$ be a natural number, and let a constraint be given by a
linear equation over the field with two elements, with a nonempty ordered tuple of
distinct variables from $V$ as its scope. For any assignment $x:V\to\mathbb{F}_2$, the
first component of its decomposition into values on the constraint scope and values
outside that scope, evaluated at position $i$, is $x$ evaluated at the variable in scope
position $i$.
\end{theorem}

\begin{theorem}[Expected satisfaction of a linear constraint]
\lean{FABL.expect_satisfiedIndicator_maxLinear}
\label{FABL.expect_satisfiedIndicator_maxLinear}
\leanok
\uses{FABL.CSPConstraint, FABL.CSPConstraint.satisfiedIndicator, FABL.expect_maxLinearPredicate, FABL.maxLinearScopeSplitEquiv, FABL.maxLinearScopeSplitEquiv_fst_apply, FABL.maxLinearTemplate, FABL.𝔽₂}
\difficulty{5}
Let $r$ be a natural number, let $V$ be a finite set, and consider a constraint on $k$
distinct variables, where $1\le k\le r$, requiring the sum of their values in
$\mathbb{F}_2$ to equal a specified right-hand side in $\mathbb{F}_2$. Under a uniformly
random assignment of elements of $\mathbb{F}_2$ to $V$, the expected satisfaction
indicator is $1/2$.
\end{theorem}

\begin{theorem}[Expected number of satisfied linear constraints]
\lean{FABL.expect_satisfiedIndicators_maxLinear}
\label{FABL.expect_satisfiedIndicators_maxLinear}
\leanok
\uses{FABL.CSPConstraint, FABL.CSPConstraint.satisfiedIndicator, FABL.expect_satisfiedIndicator_maxLinear, FABL.maxLinearTemplate, FABL.𝔽₂}
\difficulty{4}
Let $r$ be a natural number, let $V$ be a finite set, and let $\mathcal{C}$ be a finite
multiset of constraints. Each constraint uses $k$ distinct variables for some $1\le k\le
r$ and requires their sum in $\mathbb{F}_2$ to equal a specified right-hand side in
$\mathbb{F}_2$. Under a uniformly random assignment of elements of $\mathbb{F}_2$ to
$V$, the expected number of satisfied constraints in $\mathcal{C}$ is $|\mathcal{C}|/2$.
\end{theorem}

\begin{theorem}[Expect value max linear]
\lean{FABL.expect_value_maxLinear}
\statementsource{Analysis-of-Boolean-Functions-by-Ryan-ODonnell}{
  pdf-ffbe9c7b3306-p186-b001,
  pdf-ffbe9c7b3306-p198-b006
}
\label{FABL.expect_value_maxLinear}
\leanok
\uses{FABL.CSPInstance, FABL.CSPInstance.satisfiedCount, FABL.CSPInstance.value, FABL.expect_satisfiedIndicators_maxLinear, FABL.maxLinearTemplate, FABL.𝔽₂}
\difficulty{3}
% dataset-pipeline: stub -- prose pending build_informal_statements.py
Statement pending regeneration.
\end{theorem}

\begin{theorem}[Expected value of a random linear-constraint assignment]
\lean{FABL.uniformRandomizedAssignment_expectedValue_maxLinear}
\label{FABL.uniformRandomizedAssignment_expectedValue_maxLinear}
\leanok
\uses{FABL.CSPInstance, FABL.CSPInstance.value, FABL.FiniteRandomizedAssignment.sampleLaw, FABL.expect_value_maxLinear, FABL.maxLinearTemplate, FABL.pmfExpectation, FABL.pmfExpectation_uniformPMF_eq_expect, FABL.uniformRandomizedAssignment, FABL.uniformRandomizedAssignment_sampleLaw, FABL.𝔽₂}
\difficulty{3}
Let $V$ be a finite set and consider a nonempty instance of linear equations over
$\mathbb{F}_2$, with every variable occurring in some equation. Each equation involves a
positive number of pairwise distinct variables and has a right-hand side in
$\mathbb{F}_2$. Under an independent uniform assignment of values in $\mathbb{F}_2$ to
the variables, the expected fraction of satisfied equations, counting multiplicities, is
$1/2$.
\end{theorem}

\begin{theorem}[Half-Value Guarantee for Linear Constraints]
\lean{FABL.half_le_optimum_maxLinear}
\proofsource{Analysis-of-Boolean-Functions-by-Ryan-ODonnell}{pdf-ffbe9c7b3306-p188-b005}
\label{FABL.half_le_optimum_maxLinear}
\leanok
\uses{FABL.CSPInstance, FABL.CSPInstance.optimum, FABL.CSPInstance.value, FABL.CSPInstance.value_le_optimum, FABL.expect_value_maxLinear, FABL.maxLinearTemplate, FABL.𝔽₂}
\difficulty{4}
Let $r$ be a natural number and let $V$ be a finite set of variables. Consider a
nonempty finite multiset of linear constraints over $\mathbb F_2$, each involving
between one and $r+1$ distinct variables, with every variable in $V$ occurring in at
least one constraint. The maximum, over assignments of elements of $\mathbb F_2$ to the
variables, of the fraction of satisfied constraints is at least $1/2$.
\end{theorem}

\begin{theorem}[Replacement in an independent product law]
\lean{FABL.bind_replace_independentProductPMF}
\label{FABL.bind_replace_independentProductPMF}
\leanok
\uses{FABL.independentProductPMF, FABL.independentProductPMF_apply}
\difficulty{5}
Let $V$ and $D$ be finite sets, let $\mu_v$ be a probability mass function on $D$ for
each $v\in V$, and fix $v_0\in V$. First sample $a$ from $\mu_{v_0}$, then form the
independent product law using $\mu_{v_0}$ replaced by the point mass at $a$. The
resulting mixture of product laws equals the independent product law formed from the
original family $(\mu_v)_{v\in V}$.
\end{theorem}

\begin{theorem}[Replacing one coordinate in a product distribution]
\lean{FABL.pmfExpectation_replace_independentProductPMF}
\label{FABL.pmfExpectation_replace_independentProductPMF}
\leanok
\uses{FABL.bind_replace_independentProductPMF, FABL.independentProductPMF, FABL.pmfExpectation, FABL.pmfExpectation_bind}
\difficulty{1}
Let $V$ and $D$ be finite sets, let $a\in V$, let $p_v$ be a probability mass function
on $D$ for each $v\in V$, and let $F:D^V\to\mathbb{R}$. Draw a label from $p_a$, then
draw a tuple independently across coordinates using $p_v$ at every coordinate other than
$a$ and the point mass at that label at coordinate $a$. The expected value of $F$ under
this procedure equals its expectation when the tuple is drawn independently across all
coordinates using the distributions $p_v$.
\end{theorem}

\begin{theorem}[Independent product of point masses]
\lean{FABL.independentProductPMF_pure}
\label{FABL.independentProductPMF_pure}
\leanok
\uses{FABL.independentProductPMF, FABL.independentProductPMF_apply}
\difficulty{4}
Let $V$ and $D$ be finite sets, and let $a:V\to D$ be an assignment. The independent
product of the point mass at $a(v)$ in each coordinate $v$ is the point mass at the
tuple $a$.
\end{theorem}

\begin{definition}[Replacement of one marginal]
\lean{FABL.replaceMarginal}
\label{FABL.replaceMarginal}
\leanok
Given a family of probability mass functions $(\mu_v)_{v\in V}$ on a set $D$, a variable
$v_0\in V$, and a label $a\in D$, define a new family by replacing $\mu_{v_0}$ with the
point mass at $a$ and leaving every other marginal unchanged.
\end{definition}

\begin{theorem}[Arity bounded by the maximum arity]
\lean{FABL.CSPTemplate.arity_le_maxArity}
\proofsource{Analysis-of-Boolean-Functions-by-Ryan-ODonnell}{pdf-ffbe9c7b3306-p184-b004}
\label{FABL.CSPTemplate.arity_le_maxArity}
\leanok
\uses{FABL.CSPTemplate, FABL.CSPTemplate.maxArity}
\difficulty{1}
Let a finite family of predicates on a domain be given, with each predicate assigned a
natural-number arity. Every predicate's arity is at most the maximum arity in the
family.
\end{theorem}

\begin{definition}[Local table work]
\lean{FABL.CSPConstraint.localTableWork}
\label{FABL.CSPConstraint.localTableWork}
\leanok
\uses{FABL.CSPConstraint}
For a finite domain $D$, a constraint using a predicate of arity $k$ has local table
work $|D|^k+1$.
\end{definition}

\begin{definition}[Local expected satisfaction indicator]
\lean{FABL.CSPConstraint.localExpectedIndicator}
\statementsource{Analysis-of-Boolean-Functions-by-Ryan-ODonnell}{pdf-ffbe9c7b3306-p198-b006}
\label{FABL.CSPConstraint.localExpectedIndicator}
\leanok
\uses{FABL.CSPConstraint, FABL.independentProductPMF, FABL.pmfExpectation}
Let $D$ be a finite label set, let a constraint apply a specified predicate to an
ordered tuple of distinct variables, and assign to each variable a probability
distribution on $D$. The local expected satisfaction indicator is the probability that
the predicate accepts when the labels at the variables in the constraint’s scope are
drawn independently according to their assigned distributions.
\end{definition}

\begin{theorem}[Irrelevant marginal replacement]
\lean{FABL.CSPConstraint.localExpectedIndicator_replace_of_not_mem}
\label{FABL.CSPConstraint.localExpectedIndicator_replace_of_not_mem}
\leanok
\uses{FABL.CSPConstraint, FABL.CSPConstraint.localExpectedIndicator, FABL.replaceMarginal}
\difficulty{3}
Let $D$ be a finite set, let a constraint template on $D$ specify a finite family of
predicates with arities and Boolean evaluations, and let $c$ be one of its constraints
on a set of variables. Suppose a variable $v$ does not occur in the scope of $c$. For
any family of probability distributions on $D$ indexed by the variables and any label
$a\in D$, replacing the distribution at $v$ by the point mass at $a$ leaves the expected
satisfaction indicator of $c$ unchanged.
\end{theorem}

\begin{theorem}[Averaging after replacing a marginal]
\lean{FABL.CSPConstraint.average_localExpectedIndicator_replace}
\label{FABL.CSPConstraint.average_localExpectedIndicator_replace}
\leanok
\uses{FABL.CSPConstraint, FABL.CSPConstraint.localExpectedIndicator, FABL.independentProductPMF, FABL.pmfExpectation, FABL.pmfExpectation_replace_independentProductPMF, FABL.replaceMarginal}
\difficulty{5}
Let $D$ be a finite domain, let a constraint template have finitely many predicates, and
let $c$ be a constraint whose scope is an ordered tuple of distinct variables. For each
variable $v$, let $P_v$ be a probability distribution on $D$, and suppose the variable
$z$ occurs in the scope of $c$. For each label $a\in D$, replace $P_z$ by the point mass
at $a$, leaving the other distributions unchanged, and let $q(a)$ be the probability
that $c$ is satisfied when its scoped variables are drawn independently from these
distributions. Then the expectation of $q(a)$ when $a$ is drawn from $P_z$ equals the
original probability that $c$ is satisfied under independent draws from the
distributions $P_v$.
\end{theorem}

\begin{theorem}[Expected indicator under a deterministic assignment]
\lean{FABL.CSPConstraint.localExpectedIndicator_pure}
\label{FABL.CSPConstraint.localExpectedIndicator_pure}
\leanok
\uses{FABL.CSPConstraint, FABL.CSPConstraint.localExpectedIndicator, FABL.CSPConstraint.satisfiedIndicator, FABL.independentProductPMF_pure, FABL.pmfExpectation_pure}
\difficulty{3}
Let $D$ be a finite set, let a constraint template on $D$ specify a finite family of
predicates with arities and Boolean evaluations, and let $c$ be one of its constraints
on a set of variables. For any assignment $x$ of values in $D$ to the variables, if each
variable is independently sampled from the point mass at its assigned value, then the
expected satisfaction indicator of $c$ equals $1$ when $c$ is satisfied by $x$ and $0$
otherwise.
\end{theorem}

\begin{definition}[Expected number of satisfied constraints]
\lean{FABL.CSPInstance.localExpectedCount}
\label{FABL.CSPInstance.localExpectedCount}
\leanok
\uses{FABL.CSPConstraint.localExpectedIndicator, FABL.CSPInstance}
For a finite domain $D$, a constraint multiset $C$ and a family of marginals
$(\mu_v)_{v\in V}$ determine the local expected count by summing, with multiplicity, the
local expected satisfaction indicator of each constraint in $C$.
\end{definition}

\begin{definition}[Local expected value]
\lean{FABL.CSPInstance.localExpectedValue}
\statementsource{Analysis-of-Boolean-Functions-by-Ryan-ODonnell}{
  pdf-ffbe9c7b3306-p186-b001,
  pdf-ffbe9c7b3306-p198-b006
}
\label{FABL.CSPInstance.localExpectedValue}
\leanok
\uses{FABL.CSPInstance, FABL.CSPInstance.localExpectedCount}
For a constraint satisfaction instance with at least one constraint and with every
declared variable appearing in some constraint scope, and a probability distribution on
the label set for each variable, the local expected value is the sum of the constraints’
local expected satisfaction indicators, counted with multiplicity, divided by the total
number of constraints counted with multiplicity.
\end{definition}

\begin{theorem}[Pure assignment value]
\lean{FABL.CSPInstance.localExpectedValue_pure}
\label{FABL.CSPInstance.localExpectedValue_pure}
\leanok
\uses{FABL.CSPConstraint.localExpectedIndicator, FABL.CSPConstraint.localExpectedIndicator_pure, FABL.CSPConstraint.satisfiedIndicator, FABL.CSPInstance, FABL.CSPInstance.localExpectedCount, FABL.CSPInstance.localExpectedValue, FABL.CSPInstance.satisfiedCount, FABL.CSPInstance.value}
\difficulty{4}
Let a finite-domain constraint satisfaction problem have a nonempty multiset of
constraints, with every declared variable occurring in some constraint scope. For any
assignment of domain labels to the variables, the expected normalized value when each
variable is independently fixed to its assigned label equals the value of that
assignment.
\end{theorem}

\begin{definition}[Constraints containing a variable]
\lean{FABL.CSPInstance.affectedConstraints}
\label{FABL.CSPInstance.affectedConstraints}
\leanok
\uses{FABL.CSPConstraint, FABL.CSPInstance}
For a constraint multiset $C$ and a variable $v\in V$, the affected constraints are the
multiset obtained by retaining exactly those constraints whose scope contains $v$.
\end{definition}

\begin{definition}[Constraints avoiding a variable]
\lean{FABL.CSPInstance.unaffectedConstraints}
\label{FABL.CSPInstance.unaffectedConstraints}
\leanok
\uses{FABL.CSPConstraint, FABL.CSPInstance}
For a constraint multiset $C$ and a variable $v\in V$, the unaffected constraints are
the multiset obtained by retaining exactly those constraints whose scope does not
contain $v$.
\end{definition}

\begin{theorem}[Partition of constraints by variable occurrence]
\lean{FABL.CSPInstance.affectedConstraints_add_unaffectedConstraints}
\label{FABL.CSPInstance.affectedConstraints_add_unaffectedConstraints}
\leanok
\uses{FABL.CSPConstraint, FABL.CSPInstance, FABL.CSPInstance.affectedConstraints, FABL.CSPInstance.unaffectedConstraints}
\difficulty{1}
Let a constraint template on a set $D$ specify a finite family of predicates with
arities and Boolean evaluations, and let $P$ be a finite set of variables. Suppose a
constraint instance on $P$ has a nonempty multiset of constraints and every variable in
$P$ occurs in the scope of at least one constraint. For any variable $v\in P$, the
multiset of constraints whose scopes contain $v$ plus the multiset of constraints whose
scopes omit $v$ equals the full constraint multiset.
\end{theorem}

\begin{definition}[Conditional count on affected constraints]
\lean{FABL.CSPInstance.affectedConditionalCount}
\label{FABL.CSPInstance.affectedConditionalCount}
\leanok
\uses{FABL.CSPConstraint.localExpectedIndicator, FABL.CSPInstance, FABL.CSPInstance.affectedConstraints, FABL.replaceMarginal}
For a finite domain $D$, a constraint multiset $C$, marginals $(\mu_v)_{v\in V}$, a
variable $v_0\in V$, and a label $a\in D$, define the affected conditional count as the
sum, with multiplicity, of the local expected satisfaction indicators over the
constraints whose scopes contain $v_0$, using marginals with the $v_0$ marginal replaced
by the point mass at $a$.
\end{definition}

\begin{theorem}[Average conditional count of affected constraints]
\lean{FABL.CSPInstance.average_affectedConditionalCount}
\label{FABL.CSPInstance.average_affectedConditionalCount}
\leanok
\uses{FABL.CSPConstraint.average_localExpectedIndicator_replace, FABL.CSPConstraint.localExpectedIndicator, FABL.CSPInstance, FABL.CSPInstance.affectedConditionalCount, FABL.CSPInstance.affectedConstraints, FABL.pmfExpectation, FABL.pmfExpectation_add, FABL.pmfExpectation_const, FABL.replaceMarginal}
\difficulty{5}
Let $D$ be a finite set, let a constraint template on $D$ specify a finite family of
predicates with arities and Boolean evaluations, and let $P$ be a finite set of
variables. Suppose a constraint instance on $P$ has a nonempty multiset of constraints
and every variable in $P$ occurs in some constraint scope. Given a probability
distribution on $D$ for each variable and a variable $v\in P$, the expectation, with $v$
sampled from its given distribution, of the sum of expected satisfaction indicators for
constraints containing $v$, when the marginal at $v$ is replaced by a point mass at the
sampled label, equals the sum of the original expected satisfaction indicators over
those constraints.
\end{theorem}

\begin{theorem}[Averaging after marginal replacement]
\lean{FABL.CSPInstance.average_localExpectedCount_replace}
\label{FABL.CSPInstance.average_localExpectedCount_replace}
\leanok
\uses{FABL.CSPConstraint.localExpectedIndicator, FABL.CSPConstraint.localExpectedIndicator_replace_of_not_mem, FABL.CSPInstance, FABL.CSPInstance.affectedConditionalCount, FABL.CSPInstance.affectedConstraints, FABL.CSPInstance.affectedConstraints_add_unaffectedConstraints, FABL.CSPInstance.average_affectedConditionalCount, FABL.CSPInstance.localExpectedCount, FABL.CSPInstance.unaffectedConstraints, FABL.pmfExpectation, FABL.pmfExpectation_add, FABL.pmfExpectation_const, FABL.replaceMarginal}
\difficulty{5}
Let $D$ be a finite set, let a constraint template on $D$ specify a finite family of
predicates with arities and Boolean evaluations, and let $P$ be a finite set of
variables. Suppose a constraint instance on $P$ has a nonempty multiset of constraints
and every variable in $P$ occurs in some constraint scope. Given a probability
distribution on $D$ for each variable and a variable $v\in P$, the expectation over a
label sampled from the distribution at $v$ of the instance’s total expected satisfaction
count after replacing the marginal at $v$ by the point mass at that label equals the
instance’s original total expected satisfaction count.
\end{theorem}

\begin{theorem}[Averaging over a replaced marginal]
\lean{FABL.CSPInstance.average_localExpectedValue_replace}
\label{FABL.CSPInstance.average_localExpectedValue_replace}
\leanok
\uses{FABL.CSPInstance, FABL.CSPInstance.average_localExpectedCount_replace, FABL.CSPInstance.localExpectedValue, FABL.pmfExpectation, FABL.pmfExpectation_mul_const, FABL.replaceMarginal}
\difficulty{3}
Let $D$ be a finite label set, let $V$ be a set of variables, and let a constraint
satisfaction instance on $V$ have a nonempty multiset of constraints in which every
variable occurs. For any family of probability laws on $D$ indexed by $V$ and any
variable $v$, replacing the law at $v$ by the point mass at a label drawn from its
original law leaves the instance's expected fraction of satisfied constraints unchanged
on average.
\end{theorem}

\begin{theorem}[Expected value with deterministic coordinates]
\lean{FABL.CSPInstance.pmfExpectation_value_eq_localExpectedValue_of_pure_off}
\label{FABL.CSPInstance.pmfExpectation_value_eq_localExpectedValue_of_pure_off}
\leanok
\uses{FABL.CSPInstance, FABL.CSPInstance.average_localExpectedValue_replace, FABL.CSPInstance.localExpectedValue, FABL.CSPInstance.localExpectedValue_pure, FABL.CSPInstance.value, FABL.independentProductPMF, FABL.independentProductPMF_pure, FABL.pmfExpectation, FABL.pmfExpectation_pure, FABL.pmfExpectation_replace_independentProductPMF, FABL.replaceMarginal}
\difficulty{5}
Let a finite-domain constraint satisfaction problem have a nonempty multiset of
constraints, with every declared variable occurring in some constraint scope. Suppose
the variable set is finite and has decidable equality, and consider a family of
probability laws on the domain together with a finite list of variables such that every
variable outside the list has a law concentrated at some label. Then the expected value
of an assignment drawn independently from these laws equals the local expected value of
the problem.
\end{theorem}

\begin{theorem}[Expected payoff and local expected value]
\lean{FABL.CSPInstance.expectedPayoff_eq_localExpectedValue}
\label{FABL.CSPInstance.expectedPayoff_eq_localExpectedValue}
\leanok
\uses{FABL.CSPInstance, FABL.CSPInstance.localExpectedValue, FABL.CSPInstance.pmfExpectation_value_eq_localExpectedValue_of_pure_off, FABL.CSPInstance.value, FABL.FiniteRandomizedAssignment, FABL.FiniteRandomizedAssignment.expectedPayoff, FABL.independentProductPMF, FABL.pmfExpectation}
\difficulty{1}
Let $D$ and $V$ be finite sets, and let a constraint satisfaction instance on $V$ have a
nonempty multiset of constraints in which every variable occurs. If each variable is
assigned an independent random label with its specified law on $D$, then the expected
fraction of satisfied constraints equals the instance's local expected value computed
from those laws.
\end{theorem}

\begin{definition}[Best local label]
\lean{FABL.CSPInstance.bestLocalLabel}
\label{FABL.CSPInstance.bestLocalLabel}
\leanok
\uses{FABL.CSPInstance, FABL.CSPInstance.affectedConditionalCount, FABL.finiteMaximizer}
For a nonempty finite domain $D$, a constraint multiset $C$, marginals $(\mu_v)_{v\in
V}$, and a variable $v_0\in V$, define a best local label to be a label in $D$
maximizing the affected conditional count.
\end{definition}

\begin{theorem}[Best local label dominates the average]
\lean{FABL.CSPInstance.pmfExpectation_le_bestLocalLabel}
\label{FABL.CSPInstance.pmfExpectation_le_bestLocalLabel}
\leanok
\uses{FABL.CSPInstance, FABL.CSPInstance.affectedConditionalCount, FABL.CSPInstance.bestLocalLabel, FABL.le_finiteMaximizer, FABL.pmfExpectation, FABL.pmfExpectation_const, FABL.pmfExpectation_mono}
\difficulty{3}
Let a finite-domain constraint satisfaction problem have a nonempty multiset of
constraints, with every declared variable occurring in some constraint scope. Suppose
the domain is nonempty and the variable set has decidable equality. For any family of
probability laws on the domain, the expected affected-constraint count at a variable is
at most that count when the variable is fixed to the selected best local label.
\end{theorem}

\begin{theorem}[Improvement by a best local label]
\lean{FABL.CSPInstance.localExpectedValue_le_replace_bestLocalLabel}
\label{FABL.CSPInstance.localExpectedValue_le_replace_bestLocalLabel}
\leanok
\uses{FABL.CSPConstraint.localExpectedIndicator, FABL.CSPConstraint.localExpectedIndicator_replace_of_not_mem, FABL.CSPInstance, FABL.CSPInstance.affectedConstraints, FABL.CSPInstance.affectedConstraints_add_unaffectedConstraints, FABL.CSPInstance.average_affectedConditionalCount, FABL.CSPInstance.bestLocalLabel, FABL.CSPInstance.localExpectedCount, FABL.CSPInstance.localExpectedValue, FABL.CSPInstance.pmfExpectation_le_bestLocalLabel, FABL.CSPInstance.unaffectedConstraints, FABL.replaceMarginal}
\difficulty{5}
Let $D$ be a finite nonempty label set, and let a constraint satisfaction instance on a
variable set $V$ have a nonempty multiset of constraints in which every variable occurs.
For any family of probability laws on $D$ and any variable $v$, replacing the law at $v$
by the point mass at a label maximizing the expected number of satisfied constraints
involving $v$ does not decrease the instance's local expected value.
\end{theorem}

\begin{definition}[Finite randomized constraint satisfaction input]
\lean{FABL.FiniteRandomizedCSPInput}
\statementsource{Analysis-of-Boolean-Functions-by-Ryan-ODonnell}{pdf-ffbe9c7b3306-p198-b006}
\label{FABL.FiniteRandomizedCSPInput}
\leanok
\uses{FABL.CSPInstance, FABL.CSPTemplate, FABL.FiniteRandomizedAssignment}
Given a constraint template over a finite label set, a finite randomized constraint
satisfaction input consists of a natural number of variables, a constraint satisfaction
instance on the canonical variable set of that size, and a randomized assignment
specifying a marginal label distribution for each variable.
\end{definition}

\begin{definition}[Local conditional expectation state]
\lean{FABL.LocalConditionalExpectationState}
\label{FABL.LocalConditionalExpectationState}
\leanok
For types $V$ and $D$, a local conditional expectation state consists of a probability
mass function on $D$ for each variable in $V$, a selected label in $D$ for each
variable, and a natural number recording the work performed so far.
\end{definition}

\begin{definition}[Work of a local conditional step]
\lean{FABL.CSPInstance.localConditionalStepWork}
\label{FABL.CSPInstance.localConditionalStepWork}
\leanok
\uses{FABL.CSPConstraint.localTableWork, FABL.CSPInstance, FABL.CSPInstance.affectedConstraints}
Let $D$ be a finite label set, let $	ext{problem}$ be a constraint satisfaction instance
with variables $V$, and let $v\in V$. The work of the local conditional step at $v$ is
$|D|$ times the sum of the local table work over all constraints affected by $v$.
\end{definition}

\begin{definition}[Local conditional step]
\lean{FABL.CSPInstance.localConditionalStep}
\statementsource{Analysis-of-Boolean-Functions-by-Ryan-ODonnell}{pdf-ffbe9c7b3306-p199-b005}
\label{FABL.CSPInstance.localConditionalStep}
\leanok
\uses{FABL.CSPInstance, FABL.CSPInstance.bestLocalLabel, FABL.CSPInstance.localConditionalStepWork, FABL.LocalConditionalExpectationState, FABL.replaceMarginal}
Let $D$ be a finite nonempty label set, let a constraint satisfaction instance have
variable set $V$, and let a state consist of a marginal distribution on $D$ and a
selected label in $D$ for each variable, together with a natural-number work count. For
a variable $v\in V$, the local conditional step replaces its marginal by the point mass
at a label maximizing the affected conditional count, updates its selected label to that
same label, and increases the work count by the work of the step at $v$.
\end{definition}

\begin{definition}[Run of local conditional steps]
\lean{FABL.CSPInstance.runLocalConditionalSteps}
\statementsource{Analysis-of-Boolean-Functions-by-Ryan-ODonnell}{pdf-ffbe9c7b3306-p199-b005}
\label{FABL.CSPInstance.runLocalConditionalSteps}
\leanok
\uses{FABL.CSPInstance, FABL.CSPInstance.localConditionalStep, FABL.LocalConditionalExpectationState}
Let $D$ be a finite nonempty label set and let a constraint satisfaction instance have
variable set $V$. Given a finite list of variables and an initial local conditional
expectation state, the run applies the local conditional step successively to the
variables in list order and returns the resulting state; for an empty list it returns
the initial state.
\end{definition}

\begin{theorem}[Total work of conditional steps]
\lean{FABL.CSPInstance.runLocalConditionalSteps_work}
\label{FABL.CSPInstance.runLocalConditionalSteps_work}
\leanok
\uses{FABL.CSPInstance, FABL.CSPInstance.localConditionalStep, FABL.CSPInstance.localConditionalStepWork, FABL.CSPInstance.runLocalConditionalSteps, FABL.LocalConditionalExpectationState}
\difficulty{4}
Let a finite-domain constraint satisfaction problem have a nonempty multiset of
constraints, with every declared variable occurring in some constraint scope. Suppose
the domain is finite and nonempty and the variable set has decidable equality. Running
the local conditional steps through any list from an initial state increases the work by
the sum of the step-work costs for the variables in that list.
\end{theorem}

\begin{theorem}[Monotonicity of local conditional steps]
\lean{FABL.CSPInstance.localExpectedValue_le_runLocalConditionalSteps}
\proofsource{Analysis-of-Boolean-Functions-by-Ryan-ODonnell}{pdf-ffbe9c7b3306-p199-b005}
\label{FABL.CSPInstance.localExpectedValue_le_runLocalConditionalSteps}
\leanok
\uses{FABL.CSPInstance, FABL.CSPInstance.localConditionalStep, FABL.CSPInstance.localExpectedValue, FABL.CSPInstance.localExpectedValue_le_replace_bestLocalLabel, FABL.CSPInstance.runLocalConditionalSteps, FABL.LocalConditionalExpectationState}
\difficulty{4}
Let $D$ be a finite nonempty label set, let $V$ be a variable set, and let a constraint
satisfaction instance on $V$ have a nonempty multiset of constraints in which every
variable occurs. Starting from any state consisting of a family of probability laws on
$D$, a label assignment, and a natural-number work total, applying the local conditional
step successively to the variables in any finite list produces a state whose local
expected value is at least that of the initial state.
\end{theorem}

\begin{theorem}[Unvisited variables are preserved]
\lean{FABL.CSPInstance.runLocalConditionalSteps_preserves_of_not_mem}
\label{FABL.CSPInstance.runLocalConditionalSteps_preserves_of_not_mem}
\leanok
\uses{FABL.CSPInstance, FABL.CSPInstance.localConditionalStep, FABL.CSPInstance.runLocalConditionalSteps, FABL.LocalConditionalExpectationState, FABL.replaceMarginal}
\difficulty{4}
Let a finite-domain constraint satisfaction problem have a nonempty multiset of
constraints, with every declared variable occurring in some constraint scope. Suppose
the domain is nonempty and the variable set has decidable equality. After running the
local conditional steps through a list from an initial state, any variable absent from
the list has both its law and its assigned label unchanged.
\end{theorem}

\begin{theorem}[Visited variables become deterministic]
\lean{FABL.CSPInstance.runLocalConditionalSteps_law_eq_pure_assignment}
\label{FABL.CSPInstance.runLocalConditionalSteps_law_eq_pure_assignment}
\leanok
\uses{FABL.CSPInstance, FABL.CSPInstance.localConditionalStep, FABL.CSPInstance.runLocalConditionalSteps, FABL.CSPInstance.runLocalConditionalSteps_preserves_of_not_mem, FABL.LocalConditionalExpectationState, FABL.replaceMarginal}
\difficulty{5}
Let a finite-domain constraint satisfaction problem have a nonempty multiset of
constraints, with every declared variable occurring in some constraint scope. Suppose
the domain is nonempty and the variable set has decidable equality. After running the
local conditional steps through a list of variables from any initial state, each
variable in that list has a law concentrated at the label assigned to it by the
resulting state.
\end{theorem}

\begin{definition}[Expected value of a randomized CSP input]
\lean{FABL.FiniteRandomizedCSPInput.expectedValue}
\statementsource{Analysis-of-Boolean-Functions-by-Ryan-ODonnell}{pdf-ffbe9c7b3306-p198-b006}
\label{FABL.FiniteRandomizedCSPInput.expectedValue}
\leanok
\uses{FABL.CSPInstance.value, FABL.FiniteRandomizedAssignment.expectedPayoff, FABL.FiniteRandomizedCSPInput}
Given a constraint template over a finite label set and a finite randomized constraint
satisfaction input for that template, its expected value is the expected payoff of the
randomized assignment when the payoff of each assignment is the number of satisfied
constraints divided by the total number of constraints, with multiplicity counted.
\end{definition}

\begin{theorem}[Expected Value as Local Expected Value]
\lean{FABL.FiniteRandomizedCSPInput.expectedValue_eq_localExpectedValue}
\label{FABL.FiniteRandomizedCSPInput.expectedValue_eq_localExpectedValue}
\leanok
\uses{FABL.CSPInstance.expectedPayoff_eq_localExpectedValue, FABL.CSPInstance.localExpectedValue, FABL.FiniteRandomizedCSPInput, FABL.FiniteRandomizedCSPInput.expectedValue}
\difficulty{1}
Let a finite CSP template over a finite label set be given, together with a finite CSP
instance having at least one constraint and a randomized assignment that independently
samples each variable according to its prescribed marginal law. The expected fraction of
satisfied constraint occurrences equals the sum of the constraints’ local satisfaction
probabilities divided by the number of constraint occurrences.
\end{theorem}

\begin{definition}[Initial local state]
\lean{FABL.FiniteRandomizedCSPInput.initialLocalState}
\label{FABL.FiniteRandomizedCSPInput.initialLocalState}
\leanok
\uses{FABL.FiniteRandomizedCSPInput, FABL.LocalConditionalExpectationState}
Let $D$ be a finite, nonempty label set, and let an input consist of a finite randomized
constraint satisfaction problem on $n$ variables together with a marginal probability
law on $D$ for each variable. Its initial local state has those same marginal laws,
assigns every variable one fixed label from $D$, and records zero work.
\end{definition}

\begin{definition}[Local conditional expectation controller]
\lean{FABL.FiniteRandomizedCSPInput.localController}
\statementsource{Analysis-of-Boolean-Functions-by-Ryan-ODonnell}{pdf-ffbe9c7b3306-p199-b005}
\label{FABL.FiniteRandomizedCSPInput.localController}
\leanok
\uses{FABL.CSPInstance.runLocalConditionalSteps, FABL.FiniteRandomizedCSPInput, FABL.FiniteRandomizedCSPInput.initialLocalState, FABL.LocalConditionalExpectationState}
Let $D$ be a finite nonempty label set and let an input specify a finite constraint
satisfaction problem together with a probability law on $D$ for each variable. The local
controller is the state obtained by running the local conditional expectation procedure
through the canonical list of all variables, starting from the input's initial state.
\end{definition}

\begin{definition}[Total conditional expectation work]
\lean{FABL.FiniteRandomizedCSPInput.conditionalExpectationWork}
\label{FABL.FiniteRandomizedCSPInput.conditionalExpectationWork}
\leanok
\uses{FABL.CSPInstance.localConditionalStepWork, FABL.FiniteRandomizedCSPInput}
Let an input consist of a finite randomized constraint satisfaction problem with
variables indexed by $\{0,\ldots,n-1\}$ and a finite label set $D$. Its conditional
expectation work is the sum, over all variables, of the work of the local conditional
step at that variable.
\end{definition}

\begin{theorem}[Controller Work Identity]
\lean{FABL.FiniteRandomizedCSPInput.localController_work_eq_conditionalExpectationWork}
\label{FABL.FiniteRandomizedCSPInput.localController_work_eq_conditionalExpectationWork}
\leanok
\uses{FABL.CSPInstance.runLocalConditionalSteps_work, FABL.FiniteRandomizedCSPInput, FABL.FiniteRandomizedCSPInput.conditionalExpectationWork, FABL.FiniteRandomizedCSPInput.initialLocalState, FABL.FiniteRandomizedCSPInput.localController}
\difficulty{3}
Let $D$ be a finite nonempty label set, let a finite CSP template and an input on
finitely many variables be given. The work of the input’s local controller equals the
total conditional-expectation work, namely the sum, over all variables, of the work of
the corresponding local conditional step.
\end{theorem}

\begin{definition}[Conditional expectation work polynomial]
\lean{FABL.FiniteRandomizedCSPInput.conditionalExpectationWorkPolynomial}
\label{FABL.FiniteRandomizedCSPInput.conditionalExpectationWorkPolynomial}
\leanok
\uses{FABL.CSPTemplate.maxArity}
For a finite label set $D$ and a constraint template whose maximum predicate arity is
$r$, the conditional expectation work polynomial is $|D|\bigl(|D|^r+1\bigr)X^2$,
regarded as a polynomial with natural-number coefficients.
\end{definition}

\begin{theorem}[Polynomial Bound on Conditional-Expectation Work]
\lean{FABL.FiniteRandomizedCSPInput.conditionalExpectationWork_le_polynomial}
\label{FABL.FiniteRandomizedCSPInput.conditionalExpectationWork_le_polynomial}
\leanok
\uses{FABL.CSPConstraint.localTableWork, FABL.CSPInstance.affectedConstraints, FABL.CSPInstance.localConditionalStepWork, FABL.CSPTemplate.arity_le_maxArity, FABL.CSPTemplate.maxArity, FABL.FiniteRandomizedCSPInput, FABL.FiniteRandomizedCSPInput.conditionalExpectationWork, FABL.FiniteRandomizedCSPInput.conditionalExpectationWorkPolynomial}
\difficulty{5}
Let $D$ be a finite nonempty label set, let a finite CSP template over $D$ have a finite
family of predicates with specified arities, and let an input consist of a CSP instance
on $n$ variables. The total conditional-expectation work is at most
$|D|\bigl(|D|^{r}+1\bigr)(n+m)^2$, where $r$ is the maximum predicate arity and $m$ is
the number of constraint occurrences.
\end{theorem}

\begin{definition}[Conditional expectation plan]
\lean{FABL.FiniteRandomizedCSPInput.ConditionalExpectationPlan}
\label{FABL.FiniteRandomizedCSPInput.ConditionalExpectationPlan}
\leanok
\uses{FABL.FiniteRandomizedCSPInput}
For a finite randomized constraint satisfaction input with finite label set $D$, a
conditional expectation plan consists of an assignment of a label in $D$ to each
variable and a natural number recording the work computed by the controller.
\end{definition}

\begin{definition}[Conditional expectation plan construction]
\lean{FABL.FiniteRandomizedCSPInput.conditionalExpectationPlan}
\label{FABL.FiniteRandomizedCSPInput.conditionalExpectationPlan}
\leanok
\uses{FABL.FiniteRandomizedCSPInput, FABL.FiniteRandomizedCSPInput.ConditionalExpectationPlan, FABL.FiniteRandomizedCSPInput.localController}
Let $D$ be a finite, nonempty label set, and let an input be a finite randomized
constraint satisfaction problem on $n$ variables with a marginal label law for each
variable. Its conditional expectation plan is the plan whose assignment and work are
those produced by the input's local controller.
\end{definition}

\begin{theorem}[Conditional-Expectation Value Guarantee]
\lean{FABL.FiniteRandomizedCSPInput.expectedValue_le_conditionalExpectationPlan}
\proofsource{Analysis-of-Boolean-Functions-by-Ryan-ODonnell}{pdf-ffbe9c7b3306-p199-b005}
\label{FABL.FiniteRandomizedCSPInput.expectedValue_le_conditionalExpectationPlan}
\leanok
\uses{FABL.CSPInstance.localExpectedValue, FABL.CSPInstance.localExpectedValue_le_runLocalConditionalSteps, FABL.CSPInstance.localExpectedValue_pure, FABL.CSPInstance.runLocalConditionalSteps_law_eq_pure_assignment, FABL.CSPInstance.value, FABL.FiniteRandomizedCSPInput, FABL.FiniteRandomizedCSPInput.conditionalExpectationPlan, FABL.FiniteRandomizedCSPInput.expectedValue, FABL.FiniteRandomizedCSPInput.expectedValue_eq_localExpectedValue, FABL.FiniteRandomizedCSPInput.initialLocalState, FABL.FiniteRandomizedCSPInput.localController}
\difficulty{3}
Let $D$ be a finite nonempty label set, let a finite CSP template over $D$ and a finite
CSP instance with at least one constraint be given, and let the input include a
randomized assignment. The expected fraction of satisfied constraints is at most the
fraction satisfied by the assignment produced by the input’s conditional-expectation
plan.
\end{theorem}

\begin{theorem}[Polynomial Bound on Plan Work]
\lean{FABL.FiniteRandomizedCSPInput.conditionalExpectationPlan_work_le_polynomial}
\label{FABL.FiniteRandomizedCSPInput.conditionalExpectationPlan_work_le_polynomial}
\leanok
\uses{FABL.FiniteRandomizedCSPInput, FABL.FiniteRandomizedCSPInput.conditionalExpectationPlan, FABL.FiniteRandomizedCSPInput.conditionalExpectationWorkPolynomial, FABL.FiniteRandomizedCSPInput.conditionalExpectationWork_le_polynomial, FABL.FiniteRandomizedCSPInput.localController_work_eq_conditionalExpectationWork}
\difficulty{1}
Let $D$ be a finite nonempty label set, let a finite CSP template over $D$ have a finite
family of predicates with specified arities, and let an input consist of a CSP instance
on $n$ variables and a randomized assignment. The work of its conditional-expectation
plan is at most $|D|\bigl(|D|^{r}+1\bigr)(n+m)^2$, where $r$ is the maximum predicate
arity and $m$ is the number of constraint occurrences.
\end{theorem}

\begin{definition}[Packed finite assignment]
\lean{FABL.PackedFiniteAssignment}
\label{FABL.PackedFiniteAssignment}
\leanok
For a type $D$, a packed finite assignment consists of a natural number $n$ and a choice
of a label in $D$ for each of the $n$ variables indexed by $\{0,\ldots,n-1\}$.
\end{definition}

\begin{definition}[Packed assignment from a plan]
\lean{FABL.FiniteRandomizedCSPInput.ConditionalExpectationPlan.toPackedAssignment}
\label{FABL.FiniteRandomizedCSPInput.ConditionalExpectationPlan.toPackedAssignment}
\leanok
\uses{FABL.FiniteRandomizedCSPInput, FABL.FiniteRandomizedCSPInput.ConditionalExpectationPlan, FABL.PackedFiniteAssignment}
Given a conditional-expectation plan for a finite randomized constraint-satisfaction
input with label set $D$, the associated packed assignment has the same number of
variables and assigns each variable the label specified by the plan.
\end{definition}

\begin{definition}[Conditional-expectation polynomial-time adapter]
\lean{FABL.ConditionalExpectationPolytimeAdapter}
\label{FABL.ConditionalExpectationPolytimeAdapter}
\leanok
\uses{FABL.BinaryEncoding, FABL.CSPTemplate, FABL.FiniteRandomizedCSPInput, FABL.FiniteRandomizedCSPInput.ConditionalExpectationPlan.toPackedAssignment, FABL.FiniteRandomizedCSPInput.conditionalExpectationPlan, FABL.PackedFiniteAssignment, FABL.PolytimeMap}
Let $D$ be a finite, nonempty label set, let a constraint template over $D$ and an
encoded-input type be given, and fix binary encodings of the inputs and packed
assignments together with a map from encoded inputs to finite randomized
constraint-satisfaction inputs. A polynomial-time map on binary strings is an adapter
if, on every encoded input, its output is exactly the canonical binary encoding of the
packed assignment obtained from the conditional-expectation plan for the associated
randomized constraint-satisfaction input.
\end{definition}

\begin{theorem}[Correct decoding of the algorithm output]
\lean{FABL.ConditionalExpectationPolytimeAdapter.decode_algorithm}
\label{FABL.ConditionalExpectationPolytimeAdapter.decode_algorithm}
\leanok
\uses{FABL.ConditionalExpectationPolytimeAdapter, FABL.FiniteRandomizedCSPInput.ConditionalExpectationPlan.toPackedAssignment, FABL.FiniteRandomizedCSPInput.conditionalExpectationPlan}
\difficulty{1}
Let an input type and a packed-assignment type have binary encodings, each with decoding
inverse to encoding and with every successfully decoded string equal to the canonical
encoding of its decoded object. Let a conditional-expectation adapter consist of an
algorithm on binary strings and a semantic map from encoded inputs to finite randomized
CSP inputs, and suppose the algorithm maps each canonically encoded input to the
canonical encoding of the packed assignment in its conditional-expectation plan. Then,
for every input, decoding the algorithm's output gives exactly that packed assignment.
\end{theorem}

\begin{theorem}[Decoded output preserves expected value]
\lean{FABL.ConditionalExpectationPolytimeAdapter.decoded_output_value_ge_expected}
\proofsource{Analysis-of-Boolean-Functions-by-Ryan-ODonnell}{pdf-ffbe9c7b3306-p199-b005}
\label{FABL.ConditionalExpectationPolytimeAdapter.decoded_output_value_ge_expected}
\leanok
\uses{FABL.CSPInstance.value, FABL.ConditionalExpectationPolytimeAdapter, FABL.ConditionalExpectationPolytimeAdapter.decode_algorithm, FABL.FiniteRandomizedCSPInput.ConditionalExpectationPlan.toPackedAssignment, FABL.FiniteRandomizedCSPInput.conditionalExpectationPlan, FABL.FiniteRandomizedCSPInput.expectedValue, FABL.FiniteRandomizedCSPInput.expectedValue_le_conditionalExpectationPlan}
\difficulty{1}
Let $D$ be a finite nonempty label set, and consider a finite predicate template, a
finite randomized CSP input, binary input and output encodings with decoding inverse to
encoding and canonical successful decodings, and an adapter algorithm that maps each
encoded input to the encoding of the packed assignment in its conditional-expectation
plan. For every encoded input, decoding the algorithm's output gives that packed
assignment, and the value of the plan's assignment is at least the expected CSP value of
the randomized assignment.
\end{theorem}

\begin{theorem}[Correct output, value guarantee, and work bound]
\lean{FABL.ConditionalExpectationPolytimeAdapter.decoded_output_correct_and_work_le_polynomial}
\label{FABL.ConditionalExpectationPolytimeAdapter.decoded_output_correct_and_work_le_polynomial}
\leanok
\uses{FABL.CSPInstance.value, FABL.ConditionalExpectationPolytimeAdapter, FABL.ConditionalExpectationPolytimeAdapter.decoded_output_value_ge_expected, FABL.FiniteRandomizedCSPInput.ConditionalExpectationPlan.toPackedAssignment, FABL.FiniteRandomizedCSPInput.conditionalExpectationPlan, FABL.FiniteRandomizedCSPInput.conditionalExpectationPlan_work_le_polynomial, FABL.FiniteRandomizedCSPInput.conditionalExpectationWorkPolynomial, FABL.FiniteRandomizedCSPInput.expectedValue}
\difficulty{1}
Let $D$ be a finite nonempty label set, let a finite predicate template have maximum
arity $r$, and consider a finite CSP input with $n$ variables and $m$ constraints,
supplied with a randomized assignment. Suppose binary input and output encodings have
decoding inverse to encoding and every successfully decoded string is canonical, and an
adapter algorithm maps each encoded input to the encoding of the packed assignment in
its conditional-expectation plan. For every encoded input, decoding the algorithm's
output gives that packed assignment, the value of the plan's assignment is at least the
expected CSP value of the randomized assignment, and the plan's work is at most
$|D|(|D|^r+1)(n+m)^2$.
\end{theorem}
